\section{Connected local components and propagation in rank two}
\label{sec:arithmetic-connectedness}

\paragraph{The conclusion and the rings on which it holds.}
This chapter proves the original rank-two component theorem.  A finite
weight ideal which satisfies the original curve-propagation condition
and has the original arithmetic seed lies on every type-I component.
The proof constructs the local component chains, bounds the bad loci,
finds the required characteristic-prime curves, and carries the seed
to the connected closed fibre.  Its automorphic hypotheses remain the
original hypotheses; they are not consequences of connectedness.

Fix a complete coefficient DVR \(O\), with finite residue field
\(k\) of odd characteristic \(\ell\), and a frame
\(F.\mathrm{Ok}(C_0)\).  For either permitted rank \(n=2,3\),
let \(S\) be the supplied universal space.  Write
\[
 A=S.A,\quad B=S.B,\quad \mathfrak m=\mathfrak m_B,
 \quad u=S.u,\quad d=F.d=\sum_v d_v.
                                                        \tag{ACC.1}
\]
The ring \(A\) is the actual algebra generated over \(B\) by the
matrix entries.  RTC7 gives its finite generation, Noetherianity and
direct \(\mathbb Z\)-grading, with \(A_0=B\), finite \(B\)-modules
\(A_j\), and entry weight \(w_i-w_j\).  In rank two the weights
are \(1,-1\), and in rank three they are \(1,0,-1\).
Every local character lies in \(B\), and every balanced entry monomial
lies in \(B\).  These are the upstream results from the supplied
space; no \(\mathrm{FinA}\) assumption is used.

For every original local object \(L:S.\mathrm{Loc}(x)\), set
\[
 \begin{gathered}
 O'=L.O',\quad R=L.R,\quad \iota:A\longrightarrow R,
 \quad b=\iota|_B,\quad \rho=\iota_*u,\\
 \pi_L(P)=b^{-1}(P),\qquad \operatorname{im}_L(P)=\iota^{-1}(P).
 \end{gathered}
                                                        \tag{ACC.2}
\]
ALC4 and ALC11 identify this very \(R\) with the completion of the
selected localization of \(O'\otimes_O A\), and give flatness and
\(x\cap B=\mathfrak m\).  The finite unramified \(O'/O\) and
its residue map are those of \(L\).  ALC17 gives continuity, the
residual representation, and the exact contraction of all named
entry ideals.  In particular, if the lower-entry ideal
\(\mathfrak t_+=u.\mathrm{triP}\) is not contained in \(x\),
its extension to \(R\) is the unit ideal.  Write
\(\mathfrak t_-=u.\mathrm{triM}\) for the upper-entry ideal.
In rank two type I means avoiding both ideals.

We keep the distinction between the universal complete ring \(S.R\),
its finite-type entry algebra \(A\), and the local completion \(R\).
Transcendence degrees below concern \(A/\eta\), never the fraction
field of an arbitrary completed chart quotient.  The argument allows
equal-characteristic coefficient DVRs and includes \(\ell=3\).

\paragraph{Classical inputs.}
We use the height theorem, the dimension formula for finite-type
domains over a universally catenary domain, and Grothendieck
connectedness in the following precise form.  If \(T\) is complete
Noetherian local with irreducible spectrum of dimension \(r\), and
\(J\) is proper and generated up to radical by \(h\) elements, then
\(V(J)\setminus Z\) is connected whenever
\(\dim Z<r-1-h\).  This follows from the cohomological-dimension
form in Stacks Project, Lemma 52.14.1, tag
\href{https://stacks.math.columbia.edu/tag/0ECR}{0ECR}:
an \(h\)-generator \v Cech complex has no local cohomology above
degree \(h\), and an irreducible \(r\)-dimensional spectrum has
connectedness dimension \(r\).  Complete Noetherian local rings are
universally catenary by the Cohen structure theorem; see
\href{https://stacks.math.columbia.edu/tag/032C}{032C}.
The dimension formula is
\href{https://stacks.math.columbia.edu/tag/02JU}{02JU}.
These are stated classical inputs, not unproved deformation-specific
predicates.  All their applications, including the numerical thresholds
and changes of rings, are proved below.  The simplified rank-two route
was suggested by
\href{https://github.com/JoseBrox/GRH_distillation/blob/fd38313e5c7018c3c988ba20510e6b0b3bca0519/paper/shared/selmer-support-continuation-rank-two.tex}
{Brox's rank-two supplement}, at its pinned November 2026 revision.
The needed arguments are given here with their literal source scopes
and their connections to our completed chapters; no status declaration
in that project is used as a proof.

\subsection{The complete local connectedness theorem}

% BEGIN ACC1
\paragraph{ACC1. Cutting dimensions without a regular-sequence assumption.}
\label{claim:ACC1}
Let \(T\) be Noetherian local and let \(f_1,\ldots,f_h\)
belong to its maximal ideal.  Then
\[
 \dim T\le \dim T/(f_1,\ldots,f_h)+h.                    \tag{ACC.3}
\]
If \(T\) is complete and its nilradical is prime, every minimal
prime \(p\) over these elements satisfies
\[
                        \dim T/p\ge \dim T-h.            \tag{ACC.4}
\]
The same assertion holds when their ideal is replaced by any ideal
with the same radical.  Nilpotents in \(T\) are permitted.

\begin{proof}
A Noetherian local ring has finite dimension: its finitely generated
maximal ideal and the height theorem give a finite bound.  Choose a
system of parameters of \(T/(f_1,\ldots,f_h)\).  Lifting its
\(s=\dim T/(f_1,\ldots,f_h)\) elements gives an ideal generated by
\(s+h\) elements whose radical is the maximal ideal of \(T\).
The height theorem gives \(\dim T\le s+h\), proving (ACC.3).

For (ACC.4), pass to \(D=T/\sqrt{0}\).  This is a complete
Noetherian local domain with the same spectrum and dimension as
\(T\).  It is universally catenary.  A minimal prime \(\bar p\)
over the images of the \(h\) elements has height at most \(h\).
In a catenary local domain, concatenating saturated chains from
\(0\) to \(\bar p\) and from \(\bar p\) to the maximal ideal
gives
\(\dim D=\operatorname{ht}(\bar p)+\dim D/\bar p\).
This proves (ACC.4).  Passing to the reduction and replacing an
ideal by its radical preserve all the primes and quotient dimensions
in this argument.  No independence of the equations was assumed.
\end{proof}
% END ACC1

% BEGIN ACC2
\paragraph{ACC2. Full original complete-local component chain.}
\label{claim:ACC2}
Let \(T\) be complete Noetherian local with prime nilradical.
Let \(J=(f_1,\ldots,f_h)\), with all \(f_i\) in its maximal
ideal, and suppose
\[
                          \dim T\ge D+1+h.                \tag{ACC.5}
\]
Every minimal prime \(p\) over \(J\) has \(\dim T/p\ge D+1\).
Any two such primes can be connected by a finite sequence of minimal
primes over \(J\), with successive members satisfying
\[
                          \dim T/(p+q)\ge D.              \tag{ACC.6}
\]
This is the full original \(\mathrm{c1ch}\), including its general
prime-nilradical scope.  The original additional antecedent that
\(T\) is an \(\mathrm{IsCO}\)-algebra over some complete DVR is
retained when applying the source theorem, but is not needed for
this stronger statement.

\begin{proof}
ACC1 proves the component dimension assertion.  The ideal \(J\)
is proper, so it has a nonempty finite set of minimal primes.
Form the graph on this set with edges (ACC.6).  If the graph were
disconnected, partition its vertices into two nonempty classes and
let \(Y_1,Y_2\) be the unions of the corresponding irreducible
components of \(V(J)\).  Their intersection is the finite union
of the \(V(p+q)\) for pairs in different classes.  Consequently
\(Z=Y_1\cap Y_2\) has dimension less than \(D\).
Neither \(Y_i\) lies in \(Z\): its component generic points
belong to only their own minimal components.  Thus
\(V(J)\setminus Z\) is separated into the two nonempty disjoint
closed, and hence open, subsets \(Y_i\setminus Z\).

Put \(r=\dim T\).  By (ACC.5), \(r-1-h\ge D\).
The stated connectedness theorem therefore forbids this separation.
The graph is connected, so any two vertices admit a finite edge
path, with the constant path allowed.  The proof uses the whole
irreducible spectrum of \(T\), not regularity or reducedness of
\(T\).  It also covers \(h=0\) and \(D=0\); in these cases
the same dimension assertions and graph argument remain valid.
\end{proof}
% END ACC2

% BEGIN ACC3
\paragraph{ACC3. Transporting the exact component graph through a quotient.}
\label{claim:ACC3}
If \(p:T\twoheadrightarrow R\) is a surjective ring map and
the minimal primes of \(\ker p\) have dimensions at least
\(D+1\) and are connected by (ACC.6), then the minimal primes of
\(R\) have those same bounds and chains.  In the original notation,
\[
             \mathrm{Chained}(\ker p,D)
                 \Longrightarrow \mathrm{Chained}(0_R,D).
                                                               \tag{ACC.7}
\]

\begin{proof}
The inverse-image correspondence for quotient rings identifies the
prime ideals of \(R\) with the primes of \(T\) containing the
kernel.  It preserves and reflects inclusion, hence identifies the
minimal primes in question.  For every ideal \(Q\subset R\),
\(R/Q\simeq T/p^{-1}(Q)\).  Moreover
\(p^{-1}(Q+Q')=p^{-1}(Q)+p^{-1}(Q')\), since both inverse images
contain \(\ker p\).  These isomorphisms preserve dimensions of
the components and of their pairwise intersections.  Transport each
member of a finite path by the quotient correspondence.  A constant
path remains a constant path.  This proves every clause of (ACC.7).
\end{proof}
% END ACC3

% BEGIN ACC4
\paragraph{ACC4. The APD presentation gives actual chains in both ranks.}
\label{claim:ACC4}
For every original local object in (ACC.2), at either rank \(2\)
or \(3\),
\[
                  \mathrm{Chained}(0_R,nd+1).              \tag{ACC.8}
\]
In particular every minimal component has dimension at least
\(nd+2\), and the minimal-component graph with intersections of
dimension at least \(nd+1\) is connected.  This needs neither
\(F.\mathrm{Big}\) nor a nontriangularity assumption on \(x\).

\begin{proof}
APD25 gives the actual surjective \(O'\)-algebra map
\[
 O'[[X_1,\ldots,X_N]]\twoheadrightarrow R,
 \qquad N=N'+nd+1,
 \qquad \ker p=(f_1,\ldots,f_{N'}).                       \tag{ACC.9}
\]
The source is a complete Noetherian local domain of dimension
\(N+1\).  The kernel is proper, so each listed relation lies in
its maximal ideal, including any zero relations used for padding.
Apply ACC2 with \(h=N'\) and \(D=nd+1\).  Its required dimension
is exactly
\(D+1+h=N'+nd+2=N+1\).  ACC3 then proves (ACC.8).
The coefficient ring \(O'\), point \(x\), map \(\iota\) and
all ordinary characters are unchanged.  The assertion does not
provide rank three's separate even/odd \(\mathrm{EPres}\) or
its \(\mathrm{Sharp}\) record.
\end{proof}
% END ACC4

\subsection{Dimension comparisons before and after completion}

% BEGIN ACC5
\paragraph{ACC5. Full dimension comparison for the original local chart.}
\label{claim:ACC5}
For both original ranks define
\[
 t(\eta)=\operatorname{trdeg}_{\operatorname{Frac}(B/(\eta\cap B))}
                              \operatorname{Frac}(A/\eta)
\]
for a prime \(\eta\subset A\).  For every prime \(P\subset R\),
\[
 \dim R/P\le \dim B/\pi_L(P)+t(\operatorname{im}_L(P)).   \tag{ACC.10}
\]
Conversely, for every prime \(\eta\subset x\) there is a prime
\(P\subset R\) with
\[
 \operatorname{im}_L(P)=\eta,
 \qquad \dim R/P=\dim B/(\eta\cap B)+t(\eta).             \tag{ACC.11}
\]
These assertions give both clauses of the full original
\(\mathrm{loc\_dim}\), with any supplied numerical upper bound
or exact value for the transcendence degree.

\begin{proof}
Put \(E=O'\otimes_O A\), and let \(y\) be the prescribed
maximal ideal at which it is completed.  The map \(O\to O'\)
is injective: a nonzero prime kernel would be the maximal ideal of
the DVR \(O\), contradicting
\(\mathfrak m_O O'=\mathfrak m_{O'}\ne0\).  Thus \(O'\) is
torsion-free, hence flat, over \(O\).  It is finite by the original
local-object data.  Consequently \(E\) is finite and flat over
\(A\).  ALC4 identifies \(R=\widehat{E_y}\) with the actual
maps, and \(y\cap A=x\).

Let \(q\subset y\) be a prime of \(E\), with contraction
\(\eta\) in \(A\) and \(\mathfrak p\) in \(B\).
The extension \(A/\eta\subset E/q\) is integral, so their
fraction fields have the same transcendence degree over
\(\operatorname{Frac}(B/\mathfrak p)\).  The residue field at
\(y\) is finite over \(B/\mathfrak m\): it is a field of
finite type over that field and algebraic by the maximal-ideal
theorem.  The complete ring \(B\), and its domain quotient
\(B/\mathfrak p\), are universally catenary.  The finite-type
dimension formula therefore gives
\[
               \dim (E/q)_y=\dim B/\mathfrak p+t(\eta).
                                                               \tag{ACC.12}
\]

For a prime \(P\subset R\), take \(q=P\cap E\).
The ring \(R/qR\) is the completion of \((E/q)_y\).
Completion preserves the dimension of a Noetherian local ring:
a system of parameters remains primary to the maximal ideal after
completion, giving the upper bound by its number of generators,
and faithful flatness lifts a prime chain, giving the reverse bound.
Hence (ACC.12) is its dimension.  The quotient map to
\(R/P\) proves (ACC.10).

For the converse, flat going down from \(y\) gives a prime
\(q\subset y\) of \(E\) over the prescribed \(\eta\subset x\).
Choose a minimal prime \(P\) over \(qR\) whose quotient has the
largest dimension.  There are finitely many such minimal primes;
their maximum dimension equals \(\dim R/qR\).  Flat going down
for \(E_y\to R\) shows that every minimal prime over \(qR\)
contracts to \(qE_y\): a strictly larger contraction would, by
flat going down to \(qE_y\), yield a smaller prime over \(qR\),
contradicting minimality.  Thus this chosen \(P\)
contracts to \(\eta\) in \(A\), and (ACC.12) proves (ACC.11).
No algebraicity statement about \(\operatorname{Frac}(R/P)\) is
used.  Quotients and completion retain their possible nilpotents
until the explicitly chosen prime quotients are taken.
\end{proof}
% END ACC5

% BEGIN ACC6
\paragraph{ACC6. Opposite weights leave at most one parameter.}
\label{claim:ACC6}
For either original rank, every off-triangular prime \(\eta\) of
\(A\) satisfies \(t(\eta)\le1\).  This is the full original
\(\mathrm{fibre}\) theorem.  Hence every off-triangular prime
\(P\subset R\) satisfies
\[
                         \dim R/P\le\dim B/\pi_L(P)+1.
                                                               \tag{ACC.13}
\]

\begin{proof}
Work in the domain \(D=A/\eta\), and let
\(K=\operatorname{Frac}(B/(\eta\cap B))\).  Some positive-weight
entry \(u\) and negative-weight entry \(v\) survive; write their
weights as \(a>0\) and \(-b<0\).  The balanced product
\(u^b v^a\) belongs to \(K^\times\).  Thus \(u\) is algebraic
over \(K(v)\).  For an entry \(z\) of positive weight \(c\),
the product \(z^b v^c\) belongs to \(K\), so \(z\) is also
algebraic over \(K(v)\).  For an entry of negative weight
\(-c\), use \(z^a u^c\in K\).  Weight-zero entries and all
ordinary characters already lie in \(K\).  These statements allow
zero entries; their algebraic equations still hold.

Since the entries generate \(A\) over \(B\), its fraction field
modulo \(\eta\) is algebraic over the field generated by the one
element \(v\).  This proves \(t(\eta)\le1\), without requiring
\(\eta\) homogeneous.  ALC17 identifies off-triangularity of
\(P\) with that of \(\operatorname{im}_L(P)\).  Apply ACC5
to get (ACC.13).  The argument uses integer weights and never
divides by \(\ell\), by a weight, or by three.
\end{proof}
% END ACC6

\subsection{The rank-two triangular locus}

% BEGIN ACC7
\paragraph{ACC7. A polynomial description of the strict triangular fibre.}
\label{claim:ACC7}
In rank two let \(W=F.\mathrm{strS}(k)\).  Thus \(a\in W\)
is locally constant and satisfies
\[
 \begin{split}
 a(gh)&=a(g)\Psi(h)+\Psi(g)^{-1}a(h),\qquad a(g^*)=a(g),\\
 a(\operatorname{loc}_{0,v}s)
   &=\gamma_v\bigl(\Psi(\operatorname{loc}_v s)
                         -\Psi(\operatorname{loc}_v s)^{-1}\bigr)
                    \quad\text{for some }\gamma_v\in k.
 \end{split}                                                \tag{ACC.14}
\]
Then \(r=\dim_k W\le C_0\), and the actual fibre algebra is
\[
                   A/(\mathfrak m A+\mathfrak t_-)
                       \simeq k[T_1,\ldots,T_r].           \tag{ACC.15}
\]
The isomorphism keeps every diagonal and local character equal to
its prescribed residual value.

\begin{proof}
Set \(\chi=\Psi^{-2}\delta\) on \(\Gamma\), where
\(\delta\) is the sign character.  For \(a\in W\), the function
\(z(g)=a(g)/\Psi(g)\) on \(G_0\), extended with \(z(c)=0\),
is a continuous \(\chi\)-cocycle.  Indeed the first identity in
(ACC.14) is its cocycle identity on \(G_0\), and the star identity
gives \(z(cgc^{-1})=-z(g)\), the required rule for the semidirect
product.  The last identity gives the local coboundaries, with
primitive \(-\gamma_v\).  The local separation condition makes
each primitive unique.  If its strict cohomology class is zero,
then \(z(s)=(\chi(s)-1)t\); evaluation at \(c\) gives
\(-2t=0\), so \(t=0\) and \(a=0\).  This embeds \(W\)
linearly in the original strict cohomology group bounded by
\(F.\mathrm{Ok}.b1p\), proving \(r\le C_0\).

Choose a basis \(a_1,\ldots,a_r\) of \(W\).  RTC13 proves,
for every \(k\)-vector space \(M\), that all strict cocycles
with values in \(M\) are finite sums of scalar-valued strict
cocycles.  Its coefficient proof applies to this fixed basis and
gives uniquely \(a(g)=\sum_j a_j(g)m_j\).  This includes
infinite-dimensional \(M\).  The coordinate evaluations on
\(W\) span \(W^*\): a vector annihilated by every evaluation
is the zero function.  Thus finitely many evaluations express
each \(m_j\) as a \(k\)-linear combination of values of \(a\).

Apply RTC15 to \(\sum_j T_j a_j\) over the polynomial algebra
\(k[T_1,\ldots,T_r]\).  It gives a map from the left side of
(ACC.15), sending the lower entry to this universal function.
In the opposite direction, the lower entry of the quotient's own
representation is a strict cocycle: RTC10 and RTC14 apply because
\(\mathfrak m\) and the upper entries have been killed.  Write
its coefficients as \(m_j\) and send \(T_j\mapsto m_j\).
Both composites agree on every lower entry.  They agree on \(B\),
the diagonals and local characters, which are the prescribed
constants, and these together generate the quotient.  Hence the
composites are identity maps.  Equivalently, RTC16's uniqueness
checks the same assertion.  This proves (ACC.15) on the full
possibly nonreduced fibre algebra, without taking its reduction.
\end{proof}
% END ACC7

% BEGIN ACC8
\paragraph{ACC8. The actual closed coefficient fibre has dimension at most \(C_0\).}
\label{claim:ACC8}
Let \(n=2\), let \(L\) be any original local object, and assume
\(\mathfrak t_+\not\subset x\).  Then
\[
                         \dim R/\mathfrak mR\le C_0.
                                                               \tag{ACC.16}
\]
This is a strengthened rank-two version of the original fibre
dimension bound.

\begin{proof}
Some lower entry \(v\) is outside \(x\), hence \(\iota(v)\)
is a unit of \(R\).  For every upper entry \(u'\), the product
\(u'v\) is balanced and therefore belongs to \(B\).  It belongs
to \(\mathfrak m\): in the universal residual representation
all off-diagonal entries are zero.  Dividing its image by the unit
\(\iota(v)\) gives \(\iota(u')\in\mathfrak mR\).
Thus \(\rho.\mathrm{triM}\subset\mathfrak mR\), as an actual
ideal inclusion, not just after reduction.

For a prime \(P\supset\mathfrak mR\), its contraction
\(\eta\) to \(A\) contains \(\mathfrak mA+\mathfrak t_-\).
By ACC7, \(A/\eta\) is a domain quotient of a polynomial ring
in at most \(C_0\) variables over \(k\).  Its transcendence
degree is at most \(C_0\), while \(\eta\cap B=\mathfrak m\).
ACC5 gives \(\dim R/P\le C_0\).  The dimension of a Noetherian
quotient is the maximum of the dimensions of its finitely many
minimal prime quotients.  Apply this to \(R/\mathfrak mR\)
to obtain (ACC.16).  The specified residue extension in \(L\)
has not been replaced; ACC5 accounts for it and the completion.
\end{proof}
% END ACC8

% BEGIN ACC9
\paragraph{ACC9. The full finite-character detector supplied by the frame.}
\label{claim:ACC9}
For any field \(k\) of characteristic \(\ell\) and any frame
with \(F.\mathrm{Ok}(C_0)\), there are
\[
 s_1,\ldots,s_{1+C_0+d}\in G_0,\qquad
 s'_1,\ldots,s'_d\in G_0,\qquad H>0                       \tag{ACC.17}
\]
with \(\Psi(s_a)=1\), and with each \(s'_a\) the image of
an original inertia generator on which \(\Psi=1\), such that
every open-kernel character \(\xi:G_0\to C\) to a finite
abelian \(\ell\)-group satisfies both implications:
\[
 \begin{split}
 \xi(s_a)=1\ \text{for all }a
       &\Longrightarrow \xi(g)^H=1\ \text{for all }g,\\
 \xi(g^*)=\xi(g),\quad \xi(s'_a)=1\ \text{for all }a
       &\Longrightarrow \xi(g)^H=1\ \text{for all }g.
 \end{split}                                                \tag{ACC.18}
\]
This is the full original \(\mathrm{ab\_gen}\); no chosen
coefficient DVR is required.  The finiteness of \(k\) is already
one of the \(F.\mathrm{Ok}\) fields.

\begin{proof}
The space of continuous additive homomorphisms \(\Gamma\to k\)
embeds in the trivial-character \(H^1_C\) with full local
subspaces: send a homomorphism to its cocycle with zero local
primitives.  A boundary has zero global component.  The frame's
\(h1\) bound therefore makes this space have dimension at most
\(1+C_0\).  Choose \(1+C_0\) evaluations detecting it, padding
with the identity if necessary.  Call the corresponding group
elements \(t_j\).  They detect every finite abelian
\(\ell\)-group character of \(\Gamma\): a nontrivial image has
a nontrivial quotient of order \(\ell\), and composing with
\(\mathbb F_\ell\hookrightarrow k\) would give an undetected
nonzero additive character.

Choose the original \(d_v\) inertia generators at each place
from \(F.\mathrm{Ok}.gen\), and enumerate their global images
as \(s'_1,\ldots,s'_d\).  If \(\xi\) kills these generators,
the same generator property kills its restriction to every original
inertia group.  If also \(\xi(g^*)=\xi(g)\), then
\(\xi(cgc^{-1})=\xi(g)^{-1}\); it kills the conjugate inertia
groups as well.  The exact \(cl\) field supplies one nonzero
integer \(H\) killing every such unramified character.  This
proves the second implication of (ACC.18).

Write \(t_j=a_jc^{\epsilon_j}\), with \(a_j\in G_0\), and
put \(q=|k^\times|\).  As the first \(1+C_0\) elements of
\(s\), take \((a_jca_jc^{-1})^q\), and append the \(s'_a\).
All have \(\Psi=1\).  If \(\xi\) kills them, coprimality
of \(q\) and \(\ell\) shows that
\(\xi(a_j)\xi(ca_jc^{-1})=1\).  The character
\(g\mapsto\xi(g)\xi(cgc^{-1})\) is conjugation invariant,
so extends to \(\Gamma\) by sending \(c\) to one.  Its
kernel is open, since \(G_0\) and the original kernels are open.
It kills all \(t_j\), so is trivial by their detecting property.
Thus \(\xi(g^*)=\xi(g)\).  The appended inertia generators
and the second implication give the first.  The same \(H\)
works in both assertions and is chosen before every test character.
\end{proof}
% END ACC9

% BEGIN ACC10
\paragraph{ACC10. Passing finite-character tests to an adic character.}
\label{claim:ACC10}
Let \(D\) be Noetherian local of characteristic \(\ell\),
let \(G\) be a compact topological group, and let
\(\psi:G\to D^\times\) be adically continuous.  Suppose
\(q\) is prime to \(\ell\), \(\psi(g)^q\equiv1\pmod{
\mathfrak m_D}\) for all \(g\), and specified elements \(u_a\)
are sent to one.  If the original finite abelian \(\ell\)-group
generator test at these \(u_a\) forces every such character to
kill an element \(g\), and \(\psi(g)\equiv1\pmod{
\mathfrak m_D}\), then \(\psi(g)=1\).

Also, if \(D\) is a domain and a residually trivial adic character
\(\xi:G\to D^\times\) has exponent dividing one fixed \(H>0\)
in every finite adic image, then \(\xi=1\).

\begin{proof}
For each \(N\ge1\), the image of \(\psi^q\) in
\((D/\mathfrak m_D^N)^\times\) is finite: its map is locally
constant and its domain compact.  The image lies in \(1+
\mathfrak m_D/\mathfrak m_D^N\), an \(\ell\)-group on any
finite subgroup, since
\((1+z)^{\ell^a}=1+z^{\ell^a}\) and \(\ell^a\ge N\)
annihilates the second term.  The kernel is open.  Apply the
stated finite generator property to this actual finite image.
It gives \(\psi(g)^q=1\) modulo \(\mathfrak m_D^N\).
The residually trivial \(\psi(g)\) also has \(\ell\)-power
order there, so Bezout's identity gives \(\psi(g)=1\) in that
quotient.  Krull intersection gives equality in \(D\).
Neither completeness nor finiteness of the residue field was
needed in this first argument.

For the second assertion the same finite-image argument and
separatedness give \(\xi(g)^H=1\) in \(D\).  Write
\(H=\ell^a H_0\), with \(H_0\) prime to \(\ell\).
Since \(D\) is a domain,
\((\xi(g)^{H_0}-1)^{\ell^a}=0\) implies
\(\xi(g)^{H_0}=1\).  The factor
\(1+\xi(g)+\cdots+\xi(g)^{H_0-1}\) has residue \(H_0\)
and is a unit.  Its product with \(\xi(g)-1\) is zero, so
\(\xi(g)=1\).  This proves the asserted passage from bounded
finite exponents to the exact character identity.
\end{proof}
% END ACC10

% BEGIN ACC11
\paragraph{ACC11. The triangular cuts force the full coefficient ideal.}
\label{claim:ACC11}
Return to rank two and a local object with
\(\mathfrak t_+\not\subset x\).  Choose \(s_a\) from ACC9
and a uniformizer \(\varpi\) of \(O\).  Put
\[
 J=\rho.\mathrm{triM}+
       (\varpi,\ \rho(s_a)_{00}-1\ (1\le a\le1+C_0+d))
                          \subset R.                     \tag{ACC.19}
\]
Then \(\mathfrak mR\subset\sqrt J\).  All listed cuts are
in \(\mathfrak m_R\), and their number, beyond the triangular
ideal, is \(d+C_0+2\).

\begin{proof}
The coefficient map is local, and the diagonal reduction at each
\(s_a\) equals \(\Psi(s_a)=1\), so the cuts belong to
\(\mathfrak m_R\).  Let \(P\) be any prime containing \(J\)
and put \(D=R/P\).  It is a separated complete local domain of
characteristic \(\ell\), with its specified embedding of \(k\).
The triangular representation has diagonal characters \(\alpha,
\beta\).  Dividing \(\alpha\) by the embedded \(\Psi\)
gives a residually trivial adic character which kills all \(s_a\).
ACC9 applied to each finite adic image, followed by ACC10, gives
\(\alpha=\Psi\).  Polarization gives
\(\beta(g)=\alpha(g^*)=\Psi(g)^{-1}\).

The local characters also equal their prescribed values, including
their order.  To see this without choosing a new flag, use ACP4's
actual characteristic-polynomial identity.  At a local separating
element \(h\), its two diagonal values have distinct residues.
The two roots supplied by the ordered local characters are therefore
matched to these diagonal values by their residues.  For arbitrary
\(g\), compare the traces at \(g\) and \(hg\).  The resulting
two linear equations have determinant
\(\alpha(h)-\beta(h)\), a unit, and force the same ordered
matching at \(g\).  Thus both full local characters are residual.

A balanced entry monomial containing a strictly lower entry also
contains an upper entry, so its image is zero.  The other balanced
monomials, and every local character value, are the prescribed
constants.  Hence the dense generating subalgebra \(B_0\) of
\(B\) maps to its residue in \(D\).  RTC7's topology comparison
makes \(B_0\) dense for the actual maximal-ideal topology on
\(B\), and the local map \(B\to D\) is continuous.  For
\(b\in B\), its approximants have eventually the fixed residue
of \(b\); separatedness then gives the same image for \(b\).
Consequently \(\mathfrak m\) maps to zero in \(D\).
This holds for every prime \(P\supset J\), proving the radical
inclusion.  The argument has retained all character values, not
only their semisimplified characteristic polynomials.
\end{proof}
% END ACC11

% BEGIN ACC12
\paragraph{ACC12. The full rank-two triangular dimension bound.}
\label{claim:ACC12}
Under the rank-two hypotheses of ACC11,
\[
 \dim R/\rho.\mathrm{triM}\le d+2C_0+2,
 \qquad \rho.\mathrm{triP}=R.                            \tag{ACC.20}
\]
In particular the original rank-two instance of \(\mathrm{l51}\)
holds, whose numerical upper bound is \(d+2+(3C_0+5)\).
If \(F.\mathrm{Big}(C_0)\), then
\(\dim R/\rho.\mathrm{triM}<2d\).

\begin{proof}
ACC11 gives \(\dim R/J\le\dim R/\mathfrak mR\), and ACC8
bounds the latter by \(C_0\).  Apply ACC1 to the local quotient
\(R/\rho.\mathrm{triM}\) and its \(d+C_0+2\) additional
cuts.  This proves the first bound in (ACC.20).  The unit lower
entry proves the second, as in ALC17.  The displayed bound is no
larger than \(d+3C_0+7\), so gives exactly the required source
instance.  Finally \(\mathrm{Big}\) says \(d>3C_0+7\);
in particular \(d>2C_0+2\), which proves the strict inequality.
There is no assertion here that the rank-three instance of
\(\mathrm{l51}\) has been supplied.
\end{proof}
% END ACC12

% BEGIN ACC13
\paragraph{ACC13. Full original rank-two sharpness.}
\label{claim:ACC13}
For every original rank-two space and local object, if
\(F.\mathrm{Big}(C_0)\) and \(\mathfrak t_+\not\subset x\),
then
\[
 \mathrm{Chained}(0_R,2d+1),\qquad
       \rho.\mathrm{TI}(P)\quad\text{for every minimal prime }P
       \text{ of }R.                                     \tag{ACC.21}
\]
This is the full original \(\mathrm{sharp2}\).

\begin{proof}
The chain assertion is ACC4.  Let \(P\) be a minimal prime.
Its quotient has dimension at least \(2d+2\).  The lower-entry
ideal is the unit ideal, so it is not contained in \(P\).
If the upper-entry ideal lay in \(P\), the quotient map would
give
\(2d+2\le\dim R/P\le\dim R/\rho.\mathrm{triM}<2d\),
contradiction.  Thus \(P\) avoids both triangular ideals.
In rank two this is precisely type I; the extra exceptional-type
clause of the rank-three definition has false antecedent.  Every
original quantifier over \(O,F,S,x,L\) has been retained.
\end{proof}
% END ACC13

\subsection{Coincident characters and characteristic-prime curves}

% BEGIN ACC14
\paragraph{ACC14. Counting the exact ordinary-character cuts.}
\label{claim:ACC14}
This claim applies at both ranks \(n=2,3\).  Let \(I\subset B\)
satisfy the original finite-weight condition
\[
 \mathfrak m\subset\sqrt{I+\lambda},\qquad
 \lambda=\mathfrak m_O B+
  (\psi_{v,i}(g)-1:g\in I_v,\ \Psi(\operatorname{loc}_v g)=1).
                                                               \tag{ACC.22}
\]
Let \(\mathfrak p\) be any ideal of \(B\), containing \(I\),
such that for specified \(v\) and \(i\ne j\) the two ordered
characters agree on \(N_v\) modulo \(\mathfrak p\).  Then
\[
                        \dim B/\mathfrak p\le nd-d_v+2.
                                                               \tag{ACC.23}
\]
No primality or characteristic assumption on \(\mathfrak p\)
is required.  This supplies the original coefficient-ring
dimension argument \(\mathrm{dim\_B}\).

\begin{proof}
Use the original inertia generators \(u_{w,a}\), with one exceptional
index \(a_0(w)\) at each place; all nonexceptional generators lie
in \(N_w\).  Impose these equations in \(B/\mathfrak p\):
the image of \(\varpi\), every
\(\psi_{w,r}(u_{w,a})-1\) except the \(d_v\) equations for
\((w,r)=(v,j)\), and the one omitted exceptional equation
\(\psi_{v,j}(u_{v,a_0(v)})-1\).  There are exactly
\(1+(nd-d_v)+1=nd-d_v+2\) listed elements, all in the maximal
ideal.  A zero or duplicate element does not affect this bound.

In any prime quotient of these cuts, every generator equation holds.
Indeed the missing nonexceptional \(j\)-equations follow by
subtracting the \(i\)-equation and using coincidence on \(N_v\).
The quotient is a local domain of characteristic \(\ell\).
For every original local character, its \(q=|k^\times|\)-th
power is residually one.  Apply ACC10 with the exact local generator
property of \(F.\mathrm{Ok}.gen\).  It makes the character one
on every inertia element on which its prescribed residual value is
one.  Hence this prime contains \(\lambda\), as well as \(I\).
Condition (ACC.22) makes it the maximal ideal of \(B\).

The cut quotient therefore has dimension zero.  ACC1 gives (ACC.23).
If \(\mathfrak p=B\), its quotient is zero and the assertion is
automatic with the convention \(\dim0=-\infty\).  Counting the
uniformizer remains necessary in equal characteristic: the equation
\(\ell=0\) there does not impose \(\varpi=0\).
\end{proof}
% END ACC14

% BEGIN ACC15
\paragraph{ACC15. Full original coincident-prime dimension bound in both ranks.}
\label{claim:ACC15}
For either original rank, let \(P\subset R\) be prime and
off triangular.  Suppose \(I\subset\pi_L(P)\), where \(I\)
satisfies (ACC.22), and suppose its two specified distinct ordinary
characters coincide on \(N_v\) modulo \(P\).  Then
\[
                            \dim R/P+d_v\le nd+3.         \tag{ACC.24}
\]
In particular this proves both conclusions of the full original
\(\mathrm{l52}\): the bound \(nd+6\), and the bound \(nd+5\)
when \(\mathfrak m_O R\subset P\).  No extra characteristic
condition is needed for the stronger bound (ACC.24).

\begin{proof}
The full ordered character differences contract to
\(\mathfrak p=\pi_L(P)\), since the map \(b\) is the
restriction of the original \(\iota\).  Apply ACC14 to get
\(\dim B/\mathfrak p\le nd-d_v+2\).  ACC6 and ACC5 give
\(\dim R/P\le\dim B/\mathfrak p+1\).  Combining these
inequalities proves (ACC.24).  It implies the two stated weaker
source bounds, retaining the antecedent of the second rather than
changing it.  There has been no identification of the rank-two and
rank-three rings and no use of rank-two triangularity in this proof.
\end{proof}
% END ACC15

% BEGIN ACC16
\paragraph{ACC16. Every bad rank-two locus has the required strict dimension.}
\label{claim:ACC16}
Let \(n=2\), assume \(F.\mathrm{Big}(C_0)\), and suppose
\(\mathfrak t_+\not\subset x\).  Besides the triangular
bound \(\dim R/\rho.\mathrm{triM}<2d\), for every \(I\)
satisfying (ACC.22), every \(v\), and every \(i\ne j\),
\[
 \dim R/\bigl(IR+\mathfrak m_O R+C_{v,i,j}\bigr)<2d,
 \quad
 C_{v,i,j}=(\rho_{v,i}(g)-\rho_{v,j}(g):g\in N_v).
                                                               \tag{ACC.25}
\]
This is the original strict-dimension theorem at rank two.

\begin{proof}
The first assertion is ACC12.  Let \(Q\) be any prime containing
the ideal in (ACC.25).  If \(\rho.\mathrm{triM}\subset Q\),
its quotient dimension is bounded by the triangular quotient.
Otherwise it is off triangular, because the lower-entry ideal is
the unit ideal.  All hypotheses of ACC15 hold for \(Q\), so
\(\dim R/Q\le2d+3-d_v<2d\), using \(d_v>5\).
There are finitely many minimal primes over the displayed ideal;
their dimensions are integers, and the dimension of the quotient is
their maximum.  This proves the same strict inequality for the whole
quotient, including its nonreduced structure.  If that quotient is
zero the convention \(-\infty<2d\) gives the assertion directly.
\end{proof}
% END ACC16

% BEGIN ACC17
\paragraph{ACC17. Full original prime-avoidance curve lemma.}
\label{claim:ACC17}
Let \(T\) be any Noetherian local ring, let \(a\in\mathfrak m_T\),
let \(N\ge1\), and let \((J_i)\) be a finite family of ideals
with \(\dim T/J_i<N\).  For every ideal \(K\) with
\(\dim T/K\ge N+1\), there is a prime \(p\) such that
\[
 K\subset p,\quad a\in p,\quad \dim T/p=1,
                     \quad J_i\not\subset p\quad\text{for all }i.
                                                               \tag{ACC.26}
\]
This is the full original \(\mathrm{NicePrimeEx}\), without
completeness, catenarity, infinitude of a residue field, or primality
of \(K\) added.

\begin{proof}
By ACC1, \(\dim T/(K+(a))\ge N\).  Select a minimal prime
\(Q\) of \(K+(a)\) whose quotient has dimension at least \(N\).
No \(J_i\) is contained in \(Q\), since that would contradict
its dimension bound.  Choose \(b_i\in J_i\setminus Q\).
Also choose \(g\in\mathfrak m_T\setminus Q\), possible because
\(\dim T/Q\ge1\).  The product \(f=g\prod_i b_i\) belongs
to \(\mathfrak m_T\setminus Q\).  If the family is empty, use
the product one.

By the ascending chain condition choose a maximal prime \(p\)
among primes containing \(Q\) and avoiding \(f\).  In the local
domain \(D=T/p\), the image of \(f\) is nonzero and nonunit,
and every nonzero prime contains it.  We prove \(\dim D=1\).
If \(\dim D\ge2\), all minimal primes over \((f)\) have
height one, by the principal ideal theorem.  There are finitely
many, and none is the maximal ideal.  Prime avoidance supplies
\(z\) in the maximal ideal outside their union.  A minimal prime
\(q\) over \((z)\) is nonzero and has height one.  By the
defining property of \(D\), it contains \(f\), hence contains
a minimal prime over \((f)\).  Two such height-one primes in
a domain cannot be strictly nested, so \(q\) is that prime,
contradicting the choice of \(z\).  Thus \(\dim D\le1\).
The nonzero nonunit \(f\) shows its maximal ideal is nonzero,
so \(\dim D\ge1\).  This proves equality.

The chosen \(p\) contains \(K+(a)\).  If it contained any
\(J_i\), it would contain \(b_i\) and hence \(f\), a
contradiction.  This proves all clauses of (ACC.26).  Only finite
prime avoidance was used, so finite residue fields cause no issue.
\end{proof}
% END ACC17

% BEGIN ACC18
\paragraph{ACC18. Actual Nice curves above every sufficiently large rank-two ideal.}
\label{claim:ACC18}
At rank two assume the hypotheses of ACC16, and let \(I\subset B\)
satisfy (ACC.22).  For every ideal \(K\subset R\), not necessarily
prime, with
\[
                       IR\subset K,\qquad \dim R/K\ge2d+1,
\]
there is an original Nice prime \(p\supset K\).  Explicitly,
\[
 \dim R/p=1,\quad \ell\in p,\quad
 \rho.\mathrm{OffTri}(p),\quad
 \neg\rho.\mathrm{Coin}(p,v,i,j)\quad(v,\ i\ne j).
                                                               \tag{ACC.27}
\]
This is the full rank-two instance of
\(\mathrm{exists\_nice\_prime\_above\_of\_large\_dim}\).

\begin{proof}
Apply ACC17 with \(N=2d\), \(a\) the image of the actual
uniformizer \(\varpi\) of \(O\), and the following finite
family of bad ideals: \(\rho.\mathrm{triM}\) and all
\(IR+\mathfrak m_O R+C_{v,i,j}\) for \(i\ne j\).
Their strict dimension bounds are ACC16.  The frame's \(\mathrm{Big}\)
condition ensures \(N\ge1\), and the coefficient map sends
\(\varpi\) into \(\mathfrak m_R\).

The resulting prime contains \(\varpi\), hence contains
\(\mathfrak m_O R\) and \(\ell\).  It avoids the upper-entry
ideal; it also avoids the lower-entry ideal, which is the unit
ideal.  If its two local characters coincided on \(N_v\), it
would contain \(C_{v,i,j}\); together with \(IR\subset K\subset p\)
and \(\mathfrak m_OR\subset p\), this would contradict avoidance
of the corresponding bad ideal.  This proves (ACC.27) with the
literal ordered local tests.  In rank two every such prime is
type I, and ALC15 identifies its field-valued representation with
an absolutely irreducible one.  The prime itself, not a change to
a more favorable coefficient chart, is the required output.
\end{proof}
% END ACC18

% BEGIN ACC19
\paragraph{ACC19. Actual local rank-two propagation data.}
\label{claim:ACC19}
Fix the same rank-two space, an ideal \(I\) satisfying (ACC.22),
and the original condition \(P1'(S,I)\): for every original
\(x,L\), every Nice type-I prime \(p\) with
\(I\subset\pi_L(p)\), and every prime \(P\subset p\),
one has \(I\subset\pi_L(P)\).  At any chart with
\(\mathfrak t_+\not\subset x\), the literal rank-two
\(\mathrm{LocData}\) with threshold \(2d+1\), sharp ideal
zero, and empty exceptional and thick sets satisfies all ten
original \(\mathrm{Ok}\) clauses.

\begin{proof}
Good primes are those containing \(IR\), hence are upward closed.
Type-I primes avoid the two triangular ideals, hence are downward
closed.  Every Nice prime is type I.  ACC18 supplies the required
Nice prime above each good prime of dimension at least \(2d+1\).
The thick-prime and exceptional-chain clauses are vacuous.  A type-I
minimal prime is a minimal prime over zero; ACC13 supplies their
adjacency chains and shows that every minimal prime is type I.
The remaining downward-goodness clause is exactly the supplied
\(P1'\), on this chart.  These verify the ten fields individually.

For clarity, the resulting propagation also has the following
direct description.  A good prime of dimension at least \(2d+1\)
lies in a Nice prime by ACC18.  A minimal prime below it is good
by \(P1'\).  If a good minimal prime \(P\) is adjacent to
\(Q\), the ideal \(P+Q\) has dimension at least \(2d+1\)
and contains \(IR\).  ACC18 gives a Nice prime above it, and
\(P1'\) makes \(Q\) good.  ACC13's finite chain reaches every
minimal prime.  This is the actual local use of ACP14, with no
new local propagation predicate left to assume.
\end{proof}
% END ACC19

\subsection{Closed-fibre incidence and the original seed}

% BEGIN ACC20
\paragraph{ACC20. Minimal primes of the actual entry algebra are homogeneous.}
\label{claim:ACC20}
At both original ranks, every minimal prime \(\eta\) of \(A\)
is homogeneous for the actual grading.  This is the full original
\(\mathrm{min\_hom}\).

\begin{proof}
For any prime \(p\) in a directly \(\mathbb Z\)-graded ring,
let \(p^h\) consist of the elements all of whose homogeneous
components belong to \(p\).  It is a homogeneous ideal contained
in \(p\).  It is prime.  Indeed, for elements \(a,b\notin p^h\),
choose the largest indices \(r,s\) for which their components
\(a_r,b_s\) do not belong to \(p\).  Such largest indices exist
because each element has only finitely many homogeneous components.
In degree \(r+s\), the product has the term \(a_r b_s\),
which is outside \(p\).  Every other term lies in \(p\), since
its first index exceeds \(r\) or its second exceeds \(s\).
Thus the degree \(r+s\) component of \(ab\) is outside \(p\),
proving \(ab\notin p^h\).

For a minimal prime \(\eta\) of \(A\), the prime
\(\eta^h\subset\eta\) must equal \(\eta\).  Consequently
\(\eta\) is generated by its homogeneous elements, precisely
the original \(S.\mathrm{IsHom}(\eta)\) condition.  No action
of the finite residue field's multiplicative group is used to
detect degrees; the direct integer grading is essential here.
\end{proof}
% END ACC20

% BEGIN ACC21
\paragraph{ACC21. Closed points in every nonempty part of the core.}
\label{claim:ACC21}
At both original ranks put
\[
 \mathcal C=S.\mathrm{BB}
   =V(\mathfrak m A)\cap D(\mathfrak t_+),\qquad
 \mathcal Z=\{x\in\mathcal C:x\text{ is maximal in }A\}.
                                                               \tag{ACC.28}
\]
If a closed subset \(Y\subset\operatorname{Spec}A\) and an
open subset \(U\) satisfy \(Y\cap U\cap\mathcal C\ne\varnothing\),
their intersection contains a point of \(\mathcal Z\).
This is the full original \(\mathrm{jac}\).  In particular
\(\mathcal Z\) is dense in \(\mathcal C\), and every closed
set meeting \(\mathcal C\) meets \(\mathcal Z\).

\begin{proof}
The fibre \(A/\mathfrak m A\) is a finite-type algebra over
\(k\).  At a given point of the stated intersection, choose a
basic open \(D(f)\subset U\) and a lower entry \(v\) which
does not vanish.  The ring obtained by imposing the closed equations
of \(Y\) on the fibre and inverting \(fv\) is nonzero and
of finite type over \(k\).  It has a maximal ideal, and its
residue field is finite over \(k\) by the maximal-ideal theorem.
Contract that ideal to the fibre ring.  Its quotient embeds in this
finite field extension; being a \(k\)-subalgebra of a finite
dimensional field, it is itself a field.  Thus the contracted ideal
is maximal, avoids \(fv\), and contains all equations of \(Y\).
Its inverse image in \(A\) is the required maximal point.
Taking \(Y\) to be the whole spectrum proves density; taking
\(U\) to be the whole spectrum proves the final assertion.
The residue field may be a finite extension of \(k\); there is
no assertion of a \(k\)-rational point.
\end{proof}
% END ACC21

% BEGIN ACC22
\paragraph{ACC22. Every homogeneous component with a lower entry meets the core.}
\label{claim:ACC22}
At both ranks, if \(\eta\subset A\) is a homogeneous prime
and \(\mathfrak t_+\not\subset\eta\), there exists a maximal
point \(x\in\mathcal C\) with \(\eta\subset x\).
This is the full original \(\mathrm{l41}\).

\begin{proof}
Suppose, to the contrary, that
\(\mathfrak t_+\subset\sqrt{\eta+\mathfrak m A}\).
Noetherianity gives \(M\ge1\) such that
\(\mathfrak t_+^M\subset\eta+\mathfrak m A\).
Every entry has weight between \(-2\) and \(2\), and the
negative-weight entries are precisely lower entries.  A monomial
of weight at most \(-2M\) must therefore contain at least
\(M\) lower factors.  RTC7's monomial spanning description
implies
\[
                   A_j\subset\mathfrak t_+^M
                           \quad(j\le-2M).               \tag{ACC.29}
\]

Choose a lower entry \(v\notin\eta\), of weight \(e\le-1\).
Set \(j=2Me\), so \(j\le-2M\).  By (ACC.29), the finite
\(B\)-module \(Q_j=A_j/(A_j\cap\eta)\) is equal to
\(\mathfrak mQ_j\).  To see equality with this exact module,
project the inclusion \(A_j\subset\eta+\mathfrak m A\)
to degree \(j\): \(\eta\) is homogeneous, and
\(\mathfrak m\) has degree zero.  Nakayama gives \(Q_j=0\).
But \(v^{2M}\in A_j\), so \(v^{2M}\in\eta\), contradicting
primality and \(v\notin\eta\).

Thus the asserted radical containment fails.  There is a prime
containing \(\eta+\mathfrak m A\) and avoiding
\(\mathfrak t_+\).  Apply ACC21 with \(Y=V(\eta)\) to
obtain a maximal such point \(x\).  This proves the required
incidence by finite generation and the actual graded pieces, without
assuming that the whole entry algebra is finite over \(B\).
\end{proof}
% END ACC22

% BEGIN ACC23
\paragraph{ACC23. Preserving the weaker seed inclusion in rank two.}
\label{claim:ACC23}
Let \(n=2\).  Suppose \(\xi\subset A\) is a type-I prime,
\(Z\subset B\) is prime, and \(Z\subset\xi\cap B\).
There is a homogeneous type-I prime \(\eta\subset A\) with
\[
                         \eta\cap B=Z,\qquad t(\eta)=1.
                                                               \tag{ACC.30}
\]
This is the rank-two instance of the original homogeneous seed-prime
supplier.  It does not strengthen \(Z\subset\xi\cap B\) to
an equality as an input.

\begin{proof}
Choose an upper entry \(u\) and a lower entry \(v\) outside
\(\xi\).  Their balanced product \(a=uv\) belongs to \(B\)
and is outside \(\xi\cap B\), hence outside \(Z\).
In \(A_a\), both \(u\) and \(v\) are units.  Every upper
entry \(u'\) and lower entry \(v'\) satisfies
\[
                   u'=\frac{u'v}{a}u,\qquad
                   v'=(uv')u^{-1}.                        \tag{ACC.31}
\]
The coefficients in these expressions belong to \(B_a\), while
diagonal entries and all ordinary characters already belong to
\(B\).  Consequently the map
\(B_a[U,U^{-1}]\to A_a\), \(U\mapsto u\), is surjective.
It is injective as well: a finite Laurent relation is a sum of
distinct homogeneous degrees, so each coefficient times the
corresponding unit power of \(u\) vanishes separately.  Each
coefficient is therefore zero.  This proves the actual graded
isomorphism
\[
                         A_a\simeq B_a[U,U^{-1}].         \tag{ACC.32}
\]

Extend \(ZB_a\) to the Laurent ring and contract its prime
ideal to \(A\), calling the result \(\eta\).  This prime
is homogeneous, contracts to \(Z\), and avoids both \(u\)
and \(v\).  It is therefore type I.  Localizing its quotient
at \(a\) gives \((B/Z)_a[U,U^{-1}]\); hence the fraction-field
extension has transcendence degree exactly one.  All choices were
made on the same \(A,B\).  No coefficient extension or local
object has yet been chosen.
\end{proof}
% END ACC23

% BEGIN ACC24
\paragraph{ACC24. Full rank-two seed incidence with the original order of choices.}
\label{claim:ACC24}
Under the rank-two assumptions of ACC23, there exists a maximal
\(x\in\mathcal C\) such that, for every supplied original
\(L:S.\mathrm{Loc}(x)\), there exists a prime \(P\subset R\)
satisfying
\[
 \rho.\mathrm{TI}(P),\qquad \pi_L(P)=Z,\qquad
                      \dim R/P=\dim B/Z+1.               \tag{ACC.33}
\]
This is the complete rank-two instance of \(\mathrm{seedpt}\).
The point \(x\) is chosen before the universally quantified
local object and its coefficient extension.

\begin{proof}
Take the homogeneous prime \(\eta\) of ACC23.  ACC22 gives a
maximal \(x\in\mathcal C\) above it; fix this \(x\) once.
For any later \(L\) at this point, the second clause of ACC5,
with \(t(\eta)=1\), supplies a prime \(P\subset R\)
contracting to \(\eta\), with the dimension in (ACC.33).
Its coefficient contraction is exactly \(Z\) by (ACC.30).
ALC17's exact ideal comparisons make it type I because \(\eta\)
is type I.  The same reasoning applies to every allowed \(L\),
with its own specified unramified extension, without rechoosing
\(x\) or \(\eta\).  In the original seed hypothesis
\(\dim B/Z=1+2d\), this produces a good type-I prime of
dimension \(2d+2\), exceeding the required threshold \(2d+1\).
\end{proof}
% END ACC24

\subsection{Return to the actual component theorem}

% BEGIN ACC25
\paragraph{ACC25. All three actual local-to-global chart conclusions.}
\label{claim:ACC25}
At rank two retain \(F.\mathrm{Big}\), \(I\)'s original
finite-weight condition and \(P1'(S,I)\).  Define
\[
 T=D(\mathfrak t_+)\cap D(\mathfrak t_-),\qquad
 G=V(IA)\subset X=\operatorname{Spec}A.                   \tag{ACC.34}
\]
For every \(x\in\mathcal Z\), goodness propagates between
all type-I irreducible components through \(x\), and at least
one type-I component passes through \(x\).  If an original
chart at \(x\) has a good type-I prime of dimension at least
\(2d+1\), some type-I component of \(X\) is wholly good.
These are the three literal conclusions of the original
rank-two \(\mathrm{at\_pt}\) application.

\begin{proof}
ALC21 supplies an original local object at the specified maximal
point \(x\) over \(\mathfrak m\).  Its map \(A\to R\)
is flat by ALC11 and has closed-point contraction \(x\).
ACP16 therefore gives the component chart: minimal primes of
\(R\) contract to minimal primes of \(A\), and every global
component through \(x\) is obtained in this way.  ALC17 identifies
the inverse image of \(T\) with the actual local type-I locus,
and ideal contraction identifies the inverse image of \(G\)
with the local good locus.  The first set is open and the second
closed by (ACC.34).

ACC19 supplies the full local propagation conclusion at this very
chart.  ACC13 gives dimension at least \(2d+2\), hence at least
\(2d+1\), and type I for every local minimal prime.  All
hypotheses of ACP17 are now supplied.  Its first, second and third
parts give exactly the three asserted conclusions.  More concretely,
a good global component through \(x\) gives a good local minimal
prime; ACC19 makes every local minimal prime good, and contraction
gives every other global component through \(x\).  A local
minimal prime always exists and is type I, while any good local
prime above the numerical threshold starts the same argument.
This uses the canonical local map, not an abstract substitute chart.
\end{proof}
% END ACC25

% BEGIN ACC26
\paragraph{ACC26. Full original propagation theorem in rank two.}
\label{claim:ACC26}
For every original coefficient DVR, frame \(F.\mathrm{Ok}(C_0)\)
with \(F.\mathrm{Big}(C_0)\), rank-two space \(S\), and
ideal \(I\subset B\), suppose
\[
 \mathrm{FinLam}(I),\qquad P1'(S,I),\qquad
 \begin{gathered}
 \xi\subset A\text{ is a type-I prime},\quad Z\subset B\text{ is prime},\\
 I\subset Z\subset\xi\cap B,\qquad \dim B/Z=1+2d
 \end{gathered}                                             \tag{ACC.35}
\]
for some \(\xi,Z\).  Then
\[
             I\subset p\cap B
       \quad\text{for every type-I prime }p\subset A.     \tag{ACC.36}
\]
This is the full original \(\mathrm{PropagationTheoremGL2}\),
or \(\mathrm{components\_good\_rank\_two}\), at all its
original premises.

\begin{proof}
We verify the actual global propagation data on \(X\), with
type-I set \(T\) and good set \(G\) from (ACC.34), core
\(\mathcal C\), selected closed points \(\mathcal Z\), and
empty exceptional set.  Noetherianity of \(X\) is RTC7.
The inclusion \(\mathcal Z\subset\mathcal C\) is its definition.
RTC19 supplies preconnectedness of this exact core.  ACC21 shows
that every closed set meeting the core meets \(\mathcal Z\),
and that the core lies in the closure of \(\mathcal Z\).

Every relevant irreducible component meets the core.  Indeed let
its minimal prime be \(\eta\).  Relevance means that it contains
a type-I point, so \(\eta\) itself is type I by generization.
ACC20 makes \(\eta\) homogeneous, and ACC22 supplies a core
point above it.  ACC25 supplies propagation between relevant
components through a selected closed point, and shows that every
such closed point belongs to a relevant component.  These verify
the first seven original global fields individually.

For the final field use the supplied seed \(\xi,Z\), retaining
its weaker inclusion.  ACC24 chooses a closed core point before
any local chart and gives, for a chart there, a type-I prime \(P\)
with \(\pi_L(P)=Z\) and dimension \(2d+2\).  Since
\(I\subset Z\), this prime is good.  ACC25 therefore supplies
one relevant component wholly contained in \(G\).

All eight fields of the actual global data have now been proved.
ACP15 makes every relevant component good and hence gives
\(T\subset G\).  For a prime \(p\), membership in \(G\)
is exactly \(IA\subset p\), or \(I\subset p\cap B\).
This proves (ACC.36).  The proof has not constructed an automorphic
seed or \(P1'\); it has used precisely the seed and propagation
premises in the original theorem.  The same frame and marked
element are retained throughout.
\end{proof}
% END ACC26

% BEGIN ACC27
\paragraph{ACC27. The exact two-rank consumer and what remains upstream.}
\label{claim:ACC27}
In the original \(\mathrm{AbsMain}\) construction, once its
rank-two \(\mathrm{Setup}\) and \(\mathrm{ModIn}\) data have
supplied their actual \(\mathrm{FinLam}\), \(P1'\) and seed
transfers, ACC26 returns the required rank-two \(\mathrm{AllGood}\)
on that same space and ideal.  It can then be supplied to the
original rank-three component theorem.  The unconditional local
chains (ACC.8), full \(\mathrm{loc\_dim}\), \(\mathrm{fibre}\),
\(\mathrm{l52}\), \(\mathrm{min\_hom}\), \(\mathrm{jac}\)
and \(\mathrm{l41}\) are also available at rank three with
their stated original scopes.

\begin{proof}
The actual consumer uses separate maps \(\Phi_2,\Phi_3\) from
its determinant rings to the two weight-zero rings and separate
ideals \(J_2,J_3\).  Its rank-two inputs are the transferred
ideal \(\Phi_2(J_2)\), the corresponding \(\mathrm{FinLam}\),
\(P1'\), and \(\mathrm{SeedIn}\).  These match (ACC.35)
literally, so ACC26 returns
\(\mathrm{AllGood}(S_2,\Phi_2(J_2))\), the next argument
passed by \(\mathrm{AbsMain}\) to the rank-three theorem.
No equality between the two spaces, rings, coefficient extensions
or ideals is needed or asserted.  Each of the additional listed
suppliers was explicitly proved above for both ranks, so is
available on \(S_3\) independently of this rank-two return.

Rank three still requires its actual even/odd presentation and
\(\mathrm{Sharp}\) argument, its remaining triangular bound,
corner/seed incidence and the thick-chain installation.  The ordinary
APD presentation does not supply those additional parity equations.
The actual automorphic \(\mathrm{ModIn}\) and original
\(\mathrm{ModU}\), \(P\)-from-automorphy and seed producers
remain separate arithmetic work.  Once the rank-three \(\mathrm{OffFin}\)
producer is supplied, ACP11, RTC24 and AFE20 give the already
proved finite-point, coefficient-descent and finite-end returns;
UCE's downstream integral finiteness need not be reproved here.
The original \(P\) premise, Host scope \(\ell=3\), common
marking, and extra Alt witnesses \(\mathrm{WitA}\) and
\(\mathrm{WitAQ}\) are unchanged.  Thus this chapter completes
the exact rank-two component theorem and its listed subtheorems,
not a whole terminal arithmetic obligation.
\end{proof}
% END ACC27

\paragraph{Reading the new boundary.}
The substantive new rank-two construction is the component theorem
ACC26: its local chains, curve existence, homogeneous incidence and
seed localization have now been proved on the actual rings.  Only
the original finite-weight, propagation and seed assumptions remain
at its entrance.  The two-rank and terminal arithmetic joins retain
their own supplied structures.  No whole selected obligation, full
\(\mathrm{GdEnd}\), full \(\mathrm{HostEnd}\), or Alt endpoint
is claimed by the local and rank-two completion above.
